% Part: first-order-logic
% Chapter: tableaux
% Section: identity

\documentclass[../../../include/open-logic-section]{subfiles}

\begin{document}

\olfileid{fol}{tab}{ide}

\olsection{\usetoken{S}{identity}తో \usetoken{P}{tableau}}

\tetoken{తాదాత్మ్య విధేయంతో}{!!{identity}} \tetoken{టాబ్లోలకు}{!!^{tableau}s} అదనపు నిగమన నియమాలు అవసరం.
$\eq$కు నియమాలు కిందివి ($t$, $t_1$, $t_2$లు సంవృత పదాలు):

\begin{defish}
\AxiomC{}
\RightLabel{$\eq$}
\UnaryInfC{\sFmla{\True}{\eq[t][t]}}
\DisplayProof
\hfill
\AxiomC{\sFmla{\True}{\eq[t_1][t_2]}}
\noLine
\UnaryInfC{\sFmla{\True}{!A(t_1)}}
\RightLabel{$\TRule{\True}{\eq}$}
\UnaryInfC{\sFmla{\True}{!A(t_2)}}
\DisplayProof
\hfill
\AxiomC{\sFmla{\True}{\eq[t_1][t_2]}}
\noLine
\UnaryInfC{\sFmla{\False}{!A(t_1)}}
\RightLabel{$\TRule{\False}{\eq}$}
\UnaryInfC{\sFmla{\False}{!A(t_2)}}
\DisplayProof
\end{defish}
మిగతా అన్ని నియమాలకు భిన్నంగా, $\TRule{\True}{\eq}$,
$\TRule{\False}{\eq}$లకు శాఖలో ఇప్పటికే \emph{రెండు} చిహ్నిత
\tetoken{సూత్రాలు}{!!{formula}s} కనిపించాలి: $\sFmla{\True}{\eq[t_1][t_2]}$,
$\sFmla{S}{!A(t_1)}$ రెండూ.

\begin{ex}
$s$, $t$లు సంవృత పదాలైతే $\eq[s][t], !A(s) \Proves !A(t)$:
\begin{oltableau}
  [\sFmla{\False}{\formula{A}(t)}, just = \TAss
    [\sFmla{\True}{\eq[s][t]}, just = \TAss
      [\sFmla{\True}{\formula{A}(s)}, just = \TAss
        [\sFmla{\True}{\formula{A}(t)}, just={\TRule{\True}{\eq}[2, 3]}, close]
      ]
    ]
  ]
\end{oltableau}
ఇది సమానమైన వాటి ప్రతిస్థాపన సూత్రం లేదా లైబ్నిట్స్ నియమంగా
తెలిసినదై ఉండవచ్చు.

\tetoken{టాబ్లోలు}{!!^{tableau}s}, $\eq$ సమమితమైనదని కూడా నిరూపిస్తాయి; అంటే
$\eq[s_1][s_2] \Proves \eq[s_2][s_1]$:
\begin{oltableau}
  [\sFmla{\False}{\eq[s_2][s_1]}, just = \TAss
    [\sFmla{\True}{\eq[s_1][s_2]}, just = \TAss
      [\sFmla{\True}{\eq[s_1][s_1]}, just = {$\eq$}
        [\sFmla{\True}{\eq[s_2][s_1]}, just = {\TRule{\True}{\eq}[2, 3]}, close]
      ]
    ]
  ]
\end{oltableau}
ఇక్కడ వరుస $2$, $\TRule{\True}{\eq}$కు మొదటి పూర్వాపేక్ష
\tetoken{సూత్రం}{!!{formula}}~$\sFmla{\True}{\eq[s_1][s_2]}$. వరుస~$3$,
$\sFmla{\True}{!A(s_1)}$ రూపంలోని రెండవది---$!A(x)$ను
$\eq[x][s_1]$గా భావించండి; అప్పుడు $!A(s_1)$ అనేది
$\eq[s_1][s_1]$, $!A(s_2)$ అనేది $\eq[s_2][s_1]$.
\sourcecorrection{OLTETAB-009}{ఈ సమమితి ఉదాహరణలో ఆంగ్ల మూలం
వరుస~$3$ను రెండవ పూర్వాపేక్ష $\sFmla{\True}{!A(s_2)}$ రూపంలో
ఉందని చెప్పింది. కానీ వెంటనే ఇచ్చిన ప్రతిస్థాపన, ప్రదర్శిత టాబ్లో
ప్రకారం వరుస~$3$ $!A(s_1)$; $!A(s_2)$ నియమపు నిష్కర్ష. అందువల్ల
పూర్వాపేక్ష సూచికను $s_1$గా సరిచేశాం.}

$\eq$ సంక్రామకమైనదని కూడా అవి నిరూపిస్తాయి; అంటే
$\eq[s_1][s_2], \eq[s_2][s_3] \Proves \eq[s_1][s_3]$:
\begin{oltableau}
  [\sFmla{\False}{\eq[s_1][s_3]}, just = \TAss
    [\sFmla{\True}{\eq[s_1][s_2]}, just = \TAss
      [\sFmla{\True}{\eq[s_2][s_3]}, just = \TAss
        [\sFmla{\True}{\eq[s_1][s_3]}, just = {\TRule{\True}{\eq}[3, 2]}, close]
      ]
    ]
  ]
\end{oltableau}
ఈ \tetoken{టాబ్లోలో}{!!{tableau}} $\TRule{\True}{\eq}$కు మొదటి పూర్వాపేక్ష
\tetoken{సూత్రం}{!!{formula}}, వరుస~$3$లోని $\sFmla{\True}{\eq[s_2][s_3]}$
($s_2$,~$t_1$ పాత్రను; $s_3$,~$t_2$ పాత్రను పోషిస్తాయి). రెండవ
పూర్వాపేక్ష, $\sFmla{\True}{!A(s_2)}$ రూపంలో, వరుస~$2$. ఇక్కడ
$!A(x)$ను $\eq[s_1][x]$గా భావించండి; అప్పుడు $!A(s_2)$ అనేది
$\eq[s_1][s_2]$గా (అంటే వరుస~$2$గా), $!A(s_3)$ అనేది నిష్కర్షలోని
\tetoken{సూత్రం}{!!{formula}}~$\eq[s_1][s_3]$గా మారుతుంది.
\sourcecorrection{OLTETAB-010}{ఈ సంక్రామకత్వ వివరణలో ఆంగ్ల మూలం
$!A(s_2)$ను సాధారణ నియమ అక్షరాలతో $\eq[t_1][t_2]$గా గుర్తించి,
దానినే వరుస~$2$ అంది. కానీ ఈ అన్వయంలో $t_1=s_2$, $t_2=s_3$ కనుక ఆ
రూపం వరుస~$3$ను సూచిస్తుంది. $!A(x)=\eq[s_1][x]$ ప్రకారం వరుస~$2$
ఖచ్చితంగా $\eq[s_1][s_2]$; అదే రూపాన్ని స్పష్టంగా రాశాం.}
\end{ex}

\begin{prob}
కిందివాటికి సంవృత \tetoken{టాబ్లోలు}{!!{tableau}s} ఇవ్వండి:
\begin{enumerate}
\item $\sFmla{\False}{\lforall[x][\lforall[y][((x = y \land !A(x))
      \lif !A(y))]]}$
\item $\sFmla{\False}{\lexists[x][(!A(x) \land
    \lforall[y][(!A(y) \lif y = x)])]}$,\\
  $\sFmla{\True}{\lexists[x][!A(x)] \land
    \lforall[y][\lforall[z][((!A(y) \land !A(z)) \lif y = z)]]}$
\end{enumerate}
\end{prob}

\end{document}
